\documentclass[10pt]{article}
\usepackage[margin=0.9in]{geometry}
\usepackage{amsmath,amssymb,amsthm,microtype}
\usepackage[hidelinks]{hyperref}
\setlength{\emergencystretch}{3em}
\newtheorem{theorem}{Theorem}
\title{Noncontracting Tails of Unbounded Collatz Orbits:\\An Exact Reduction in Lean}
\author{Collatz proof project}
\date{September 5, 2026}
\begin{document}
\maketitle
\begin{abstract}
Using the repository's reciprocal-summability theorem, we prove that every
unbounded positive natural Collatz orbit has arbitrarily late tails whose
multiplicative coefficients never contract. Conversely, every positive
seed with no coefficient contraction has an unbounded orbit. Thus
nondivergence is equivalent to the assertion that every positive seed
has some coefficient contraction. We also prove backward transport and
one residue restriction on a least noncontracting seed. Neither side of
the global equivalence is established; nontrivial cycles remain a separate
problem. No mathematical priority claim is made.
\end{abstract}

\section{Definitions and the tail reduction}
Use the shortcut map $T(n)=n/2$ for even $n$ and $(3n+1)/2$ for odd $n$.
Let $x_t=T^t(N)$ and let $j_t(N)$ count odd states before time $t$. Write
\[
a_t(N)=\frac{2^t}{3^{j_t(N)}}.
\]
A seed is \emph{noncontracting} when $2^t\le3^{j_t(N)}$ for every $t$,
or equivalently $a_t(N)\le1$ for every $t$. This is an infinite condition;
a finite prefix check does not establish it.
\begin{theorem}\label{thm:tail}
For every $N>0$,
\[
\bigl(\forall B\;\exists t\;B<T^t(N)\bigr)
\quad\Longleftrightarrow\quad
\exists s\;\forall t\;2^t\le3^{j_t(T^s(N))}.
\]
On an unbounded orbit, such indices $s$ occur beyond every prescribed time.
This does not assert that every sufficiently late index has the property.
\end{theorem}

\subsection*{From unboundedness to a noncontracting tail}
The earlier module \texttt{OrbitSummability.lean} proves $a_t(N)\to0$
for every unbounded positive orbit. Since $a_0(N)=1$, all sufficiently
late values are less than one. A maximum on the remaining finite initial
set is therefore a global maximum. Choose an index $s$ attaining it.
The odd-count cocycle gives the exact identity
\[
a_{s+t}(N)=a_s(N)\,a_t(T^s(N)).
\]
The left side is at most $a_s(N)>0$, hence the shifted factor is at most
one. This proves the forward implication.

The argument can start from any shifted unbounded orbit. Bounded ideal
correction propagates to shifts and forces escape there, so each shift
is unbounded. Applying the same maximum argument after any prescribed
index gives arbitrarily late noncontracting tails.

\subsection*{From a noncontracting seed to unboundedness}
Suppose a positive noncontracting seed had $x_t\le B$ for every $t$.
The exact correction identity
\[
2^t x_t=3^{j_t}N+d_t
\]
would imply $d_t\le2^t x_t\le3^{j_t}B$. Therefore its ideal corrections
$c_t=d_t/3^{j_t}$ would be bounded above by $B$. The earlier reciprocal
theorem proves that bounded ideal correction forces escape to infinity,
contradicting the assumed bound on $x_t$. Thus a positive noncontracting
seed is unbounded. If it occurs on another orbit, that orbit is unbounded
as well, completing Theorem~\ref{thm:tail}.

\section{The precise global obligation}
\begin{theorem}
The following propositions are equivalent:
\begin{enumerate}
\item Every positive natural orbit is bounded above.
\item For every $N>0$, there exists $t$ such that $3^{j_t(N)}<2^t$.
\end{enumerate}
\end{theorem}
If (1) holds, the converse in Theorem~\ref{thm:tail} excludes a
noncontracting seed, proving (2). If (2) holds and an unbounded orbit
existed, its positive noncontracting tail would contradict (2).

This equivalence concerns nondivergence. It does not rule out nontrivial
cycles or prove that bounded orbits reach one. The earlier first-contraction
certificate supplies actual descent when a coefficient contraction occurs
within 65 steps, but it supplies neither a universal contraction time nor
the all-time assertion (2).

\section{Backward transport and a residue restriction}
A finite prefix of length $k$ with $2^t\le3^{j_t(N)}$ for every $t\le k$,
followed by a noncontracting tail at $T^k(N)$, is noncontracting at all
times. For $t=k+r$, multiply the inequalities for the prefix and tail,
using $j_{k+r}(N)=j_k(N)+j_r(T^k(N))$. This transport rule is proved
for arbitrary prefix lengths.

In particular, an odd predecessor of a noncontracting seed is
noncontracting: its first coefficient is $3/2$, and the remaining
coefficients are that factor times the tail coefficients.
A noncontracting seed must itself be odd, since an even first step would
have coefficient $1/2$.
\begin{theorem}
If $N$ is noncontracting and $N\equiv2\pmod3$, then
\[
M=\frac{2N-1}{3}
\]
is a smaller positive odd noncontracting seed. Consequently a least
positive noncontracting seed cannot be $2\pmod3$.
\end{theorem}
Indeed $N$ is odd and at least five, so $M$ is an odd positive integer,
$M<N$, and $3M+1=2N$. Thus $T(M)=N$ and odd backward transport applies.
This excludes the residue only for a least such seed. It does not exclude
all noncontracting seeds of that residue, and it leaves residues zero
and one modulo three unresolved.

\section{Verification and remaining work}
{\raggedright
Source: \texttt{lean/Collatz/NoncontractingTail.lean}. The central
declarations are \texttt{unbounded\_iff\_noncontracting\_tail},
\texttt{unbounded\_has\_arbitrarily\_late\_noncontracting\_tail}, and
\texttt{all\_orbits\_bounded\_iff\_all\_contract}. Backward transport
is \texttt{neverContracts\_of\_prefix}; the least-seed restriction is
\texttt{least\_noncontracting\_not\_two\_mod\_three}.
The module build and selected 139-declaration axiom audit pass using only
standard foundations. No new axiom or \texttt{sorry} is used.
Build with \texttt{lake build Collatz.NoncontractingTail} from
\texttt{lean/}.\par}

The full proof still needs an exclusion of every positive noncontracting
seed, and an exclusion of nontrivial cycles. The reduction proves that the
noncontracting class covers the escape problem; it does not prove that the
class is empty. The Collatz conjecture remains unresolved.
\end{document}
